 Capítulo  3
 Tabelas  de  páginas
 Tabelas  de  páginas  são  o  mecanismo  mais  popular  pelo  qual  o  sistema  operacional  fornece  a  cada  processo  seu  próprio  
espaço  de  endereço  privado  e  memória.  Tabelas  de  páginas  determinam  o  significado  dos  endereços  de  memória  e  
quais  partes  da  memória  física  podem  ser  acessadas.  Elas  permitem  que  o  xv6  isole  os  espaços  de  endereço  de  diferentes  
processos  e  os  multiplexe  em  uma  única  memória  física.  Tabelas  de  páginas  são  um  design  popular  porque  fornecem  
um  nível  de  indireção  que  permite  aos  sistemas  operacionais  realizarem  muitos  truques.  O  xv6  realiza  alguns  truques:  
mapeia  a  mesma  memória  (uma  página  de  trampolim)  em  vários  espaços  de  endereço  e  protege  as  pilhas  do  kernel  e  do  
usuário  com  uma  página  não  mapeada.  O  restante  deste  capítulo  explica  as  tabelas  de  páginas  que  o  hardware  RISC-V  
fornece  e  como  o  xv6  as  utiliza.
 3.1  Hardware  de  paginação
 Vale  lembrar  que  as  instruções  RISC-V  (tanto  de  usuário  quanto  de  kernel)  manipulam  endereços  virtuais.  A  RAM  da  
máquina,  ou  memória  física,  é  indexada  com  endereços  físicos.  O  hardware  da  tabela  de  páginas  RISC-V  conecta  esses  
dois  tipos  de  endereços,  mapeando  cada  endereço  virtual  a  um  endereço  físico.
 27
 O  Xv6  roda  em  Sv39  RISC-V,  o  que  significa  que  apenas  os  39  bits  inferiores  de  um  endereço  virtual  de  64  bits  são  
usados;  os  25  bits  superiores  não  são  usados.  Nesta  configuração  Sv39,  uma  tabela  de  páginas  RISC-V  é  logicamente  
uma  matriz  de  2  (134.217.728)  entradas  de  tabela  de  páginas  (PTEs).  Cada  PTE  contém  um  número  de  página  física  
(PPN)  de  44  bits  e  alguns  sinalizadores.  O  hardware  de  paginação  traduz  um  endereço  virtual  usando  os  27  bits  superiores  
dos  39  bits  para  indexar  na  tabela  de  páginas  para  encontrar  um  PTE,  criando  um  endereço  físico  de  56  bits  cujos  44  bits  
superiores  vêm  do  PPN  no  PTE  e  cujos  12  bits  inferiores  são  copiados  do  endereço  virtual  original.  A  Figura  3.1  mostra  
esse  processo  com  uma  visão  lógica  da  tabela  de  páginas  como  uma  matriz  simples  de  PTEs  (veja  a  Figura  3.2  para  uma  
história  mais  completa).  Uma  tabela  de  páginas  fornece  ao  sistema  operacional  controle  sobre  as  traduções  de  endereços  
virtuais  para  físicos  na  granularidade  de  blocos  alinhados  de  4096  (212 )  bytes.  Esse  bloco  é  chamado  de  página.
 39
 No  Sv39  RISC-V,  os  25  bits  superiores  de  um  endereço  virtual  não  são  usados  para  tradução.  O  endereço  físico  
também  tem  espaço  para  crescimento:  no  formato  PTE,  há  espaço  para  que  o  número  de  páginas  físicas  aumente  em  
mais  10  bits.  Os  projetistas  do  RISC-V  escolheram  esses  números  com  base  em  previsões  tecnológicas.  2  bytes  equivalem  
a  512  GB,  o  que  deve  ser  espaço  de  endereço  suficiente  para  aplicativos  que  executam
 31
Machine Translated by Google
 Endereço  virtual
 25
 64
 EXT
 27
 Índice
 1
 0
 12
 Desvio
 44
 10
 Bandeiras  PPN
 Tabela  de  páginas
 2^27
 56
 44
 Endereço  físico
 12
 Figura  3.1:  Endereços  virtuais  e  físicos  RISC-V,  com  uma  tabela  de  páginas  lógicas  simplificada.
 56
 em  computadores  RISC-V.  2  é  espaço  de  memória  física  suficiente  para,  em  um  futuro  próximo,  acomodar  muitos  dispositivos  
de  E/S  e  chips  DRAM.  Se  mais  for  necessário,  os  projetistas  do  RISC-V  definiram  Sv48  com  endereços  virtuais  de  48  bits  [3].
 Como  mostra  a  Figura  3.2,  uma  CPU  RISC-V  traduz  um  endereço  virtual  em  um  físico  em  três  etapas.
 Uma  tabela  de  páginas  é  armazenada  na  memória  física  como  uma  árvore  de  três  níveis.  A  raiz  da  árvore  é  uma  página  de  
tabela  de  páginas  de  4.096  bytes  que  contém  512  PTEs,  que  contêm  os  endereços  físicos  das  páginas  da  tabela  de  páginas  no  
próximo  nível  da  árvore.  Cada  uma  dessas  páginas  contém  512  PTEs  para  o  nível  final  da  árvore.
 O  hardware  de  paginação  usa  os  9  bits  superiores  dos  27  bits  para  selecionar  uma  PTE  na  página  da  tabela  de  páginas  raiz,  os  
9  bits  do  meio  para  selecionar  uma  PTE  em  uma  página  da  tabela  de  páginas  no  próximo  nível  da  árvore  e  os  9  bits  inferiores  
para  selecionar  a  PTE  final.  (No  Sv48  RISC-V,  uma  tabela  de  páginas  tem  quatro  níveis,  e  os  bits  39  a  47  de  um  índice  de  
endereço  virtual  no  nível  superior.)
 Se  qualquer  um  dos  três  PTEs  necessários  para  traduzir  um  endereço  não  estiver  presente,  o  hardware  de  paginação
 gera  uma  exceção  de  falha  de  página,  deixando  para  o  kernel  lidar  com  a  exceção  (veja  Capítulo  4).
 A  estrutura  de  três  níveis  da  Figura  3.2  permite  uma  maneira  eficiente  de  gravar  PTEs  em  termos  de  memória,  em  
comparação  com  o  design  de  nível  único  da  Figura  3.1.  No  caso  comum  em  que  grandes  intervalos  de  endereços  virtuais  não  
têm  mapeamentos,  a  estrutura  de  três  níveis  pode  omitir  diretórios  de  páginas  inteiros.  Por  exemplo,  se  um  aplicativo  usa  
apenas  algumas  páginas  começando  no  endereço  zero,  as  entradas  de  1  a  511  do  diretório  de  páginas  de  nível  superior  são  
inválidas  e  o  kernel  não  precisa  alocar  páginas  para  os  511  diretórios  de  páginas  intermediárias.  Além  disso,  o  kernel  também  
não  precisa  alocar  páginas  para  os  diretórios  de  páginas  de  nível  inferior  para  esses  511  diretórios  de  páginas  intermediárias.  
Portanto,  neste  exemplo,  o  design  de  três  níveis  economiza  511  páginas  para  diretórios  de  páginas  intermediárias  e  511  ×  512  
páginas  para  diretórios  de  páginas  de  nível  inferior.
 Embora  uma  CPU  percorra  a  estrutura  de  três  níveis  no  hardware  como  parte  da  execução  de  uma  instrução  de  
carregamento  ou  armazenamento,  uma  desvantagem  potencial  dos  três  níveis  é  que  a  CPU  deve  carregar  três  PTEs  da  
memória  para  executar  a  tradução  do  endereço  virtual  na  instrução  de  carregamento/armazenamento  para  um  endereço  físico.
 Para  evitar  o  custo  de  carregamento  de  PTEs  da  memória  física,  uma  CPU  RISC-V  armazena  em  cache  as  entradas  da  tabela  
de  páginas  em  um  Translation  Look-aside  Buffer  (TLB).
 32
Machine Translated by Google
 Endereço  virtual
 9
 EXT
 44
 511
 10
 Bandeiras  PPN
 1
 511
 L2
 44
 L1
 10
 Bandeiras  PPN
 0
 satp
 63
 Reservado
 1
 Diretório  de  páginas
 53
 0
 Diretório  de  páginas
 8910  07
 Número  da  página  física
 Endereço  físico
 9
 9
 12
 Deslocamento  L0
 44
 PPN
 12
 Desvio
 44
 511
 10
 Bandeiras  PPN
 1
 0
 6  5  4  3  
Você
 RSW
 DUXR
 G
 UM
 Diretório  de  páginas
 2 1
 C
 V
 V
 R
 C
 X
 Você-  Válido-  Legível-  Gravável-  Executável-  Usuário-G  Global
 A  -  Acessado
 UM
 D-  Acessado-  Sujo  (0  no  diretório  de  páginas)
 Reservado  para  software  de  supervisão
 Figura  3.2:  Detalhes  da  tradução  de  endereços  RISC-V.
 Cada  PTE  contém  bits  de  sinalização  que  informam  ao  hardware  de  paginação  como  o  endereço  virtual  associado  
pode  ser  usado.  PTE_V  indica  se  a  PTE  está  presente:  se  não  estiver  definida,  uma  referência  à  página  causa  uma  
exceção  (ou  seja,  não  é  permitida).  PTE_R  controla  se  as  instruções  podem  ler  na  página.  PTE_W  controla  se  as  
instruções  podem  gravar  na  página.  PTE_X  controla  se  a  CPU  pode  interpretar  o  conteúdo  da  página  como  instruções  
e  executá-las.
 PTE_U  controla  se  instruções  em  modo  usuário  têm  permissão  para  acessar  a  página;  se  PTE_U  não  estiver  definido,  
o  PTE  só  poderá  ser  usado  em  modo  supervisor.  A  Figura  3.2  mostra  como  tudo  funciona.  Os  sinalizadores  e  todas  as  
outras  estruturas  relacionadas  ao  hardware  da  página  são  definidos  em  (kernel/riscv.h).
 Para  instruir  uma  CPU  a  usar  uma  tabela  de  páginas,  o  kernel  precisa  gravar  o  endereço  físico  da  página  raiz  da  
tabela  de  páginas  no  registrador  satp.  Uma  CPU  traduzirá  todos  os  endereços  gerados  por  instruções  subsequentes  
usando  a  tabela  de  páginas  apontada  por  seu  próprio  satp.  Cada  CPU  possui  seu  próprio  satp,  permitindo  que  
diferentes  CPUs  executem  processos  diferentes,  cada  um  com  um  espaço  de  endereço  privado  descrito  por  sua  
própria  tabela  de  páginas.
 Normalmente,  um  kernel  mapeia  toda  a  memória  física  em  sua  tabela  de  páginas  para  que  possa  ler  e  escrever  
em  qualquer  local  da  memória  física  usando  instruções  de  carregamento/armazenamento.  Como  os  diretórios  de  
páginas  estão  na  memória  física,  o  kernel  pode  programar  o  conteúdo  de  uma  PTE  em  um  diretório  de  páginas  
escrevendo  no  endereço  virtual  da  PTE  usando  uma  instrução  de  armazenamento  padrão.
 Algumas  notas  sobre  termos.  Memória  física  refere-se  às  células  de  armazenamento  na  DRAM.  Um  byte  de  
memória  física  possui  um  endereço,  chamado  endereço  físico.  As  instruções  usam  apenas  endereços  virtuais,  que  o  
hardware  de  paginação  traduz  em  endereços  físicos  e,  em  seguida,  envia  ao  hardware  DRAM  para  leitura.
 33
Machine Translated by Google
 ou  armazenamento  de  gravação.  Ao  contrário  da  memória  física  e  dos  endereços  virtuais,  a  memória  virtual  não  é  um  espaço  físico
 objeto,  mas  se  refere  à  coleção  de  abstrações  e  mecanismos  que  o  kernel  fornece  para  gerenciar
 memória  física  e  endereços  virtuais.
 Endereços  virtuais
 MAXVA
 PHYSTOP  
(0x88000000)
 BASE  DO  
KERN  (0x80000000)
 0x10001000
 0x10000000
 0x0C000000
 0
 Trampolim
 Página  de  guarda
 Kstack  0
 Página  de  guarda
 Kstack  1
 ...
 Memória  livre
 Dados  do  kernel
 Texto  do  kernel
 Disco  VIRTIO
 UART0
 PLIC
 Endereços  físicos
 2^56-1
 Receita  médica--
RW---
RW
RW
RW
Receita  médica
 RW
RW
RW
0x02000000
 0x1000
 0
 Não  utilizado
 Memória  física  (RAM)
 Dispositivos  
de  E/S  não  utilizados  e  outros
 Disco  VIRTIO
 UART0
 PLIC
 CLINT
 Não  utilizado
 ROM  de  inicialização
 Não  utilizado
 Figura  3.3:  À  esquerda,  o  espaço  de  endereço  do  kernel  do  xv6.  RWX  refere-se  a  leitura,  gravação  e  execução  de  PTE
 permissões.  À  direita,  o  espaço  de  endereço  físico  RISC-V  que  o  xv6  espera  ver.
 3.2  Espaço  de  endereço  do  kernel
 O  Xv6  mantém  uma  tabela  de  páginas  por  processo,  descrevendo  o  espaço  de  endereço  do  usuário  de  cada  processo,  além  de  
uma  tabela  de  página  única  que  descreve  o  espaço  de  endereço  do  kernel.  O  kernel  configura  o  layout  do  seu  espaço  de  endereço  
para  obter  acesso  à  memória  física  e  a  vários  recursos  de  hardware  em  intervalos  previsíveis.
 Endereços  virtuais.  A  Figura  3.3  mostra  como  este  layout  mapeia  endereços  virtuais  do  kernel  para  endereços  físicos.  O  arquivo  
(kernel/memlayout.h)  declara  as  constantes  para  o  layout  de  memória  do  kernel  do  xv6.
 34
Machine Translated by Google
 O  QEMU  simula  um  computador  que  inclui  RAM  (memória  física)  começando  no  endereço  físico  0x80000000  e  
continuando  até  pelo  menos  0x88000000,  que  o  xv6  chama  de  PHYSTOP.
 A  simulação  do  QEMU  também  inclui  dispositivos  de  E/S,  como  uma  interface  de  disco.  O  QEMU  expõe  as  interfaces  do  
dispositivo  ao  software  como  registradores  de  controle  mapeados  na  memória,  localizados  abaixo  de  0x80000000  no  
espaço  de  endereço  físico.  O  kernel  pode  interagir  com  os  dispositivos  lendo/escrevendo  esses  endereços  físicos  especiais;  
tais  leituras  e  gravações  se  comunicam  com  o  hardware  do  dispositivo  em  vez  da  RAM.  O  Capítulo  4  explica  como  o  xv6  
interage  com  os  dispositivos.
 O  kernel  acessa  a  RAM  e  os  registradores  de  dispositivos  mapeados  na  memória  usando  "mapeamento  direto",  ou  
seja,  mapeando  os  recursos  em  endereços  virtuais  iguais  ao  endereço  físico.  Por  exemplo,  o  próprio  kernel  está  localizado  
em  KERNBASE=0x80000000  tanto  no  espaço  de  endereço  virtual  quanto  na  memória  física.  O  mapeamento  direto  
simplifica  o  código  do  kernel  que  lê  ou  grava  na  memória  física.  Por  exemplo,  quando  o  fork  aloca  memória  do  usuário  
para  o  processo  filho,  o  alocador  retorna  o  endereço  físico  dessa  memória;  o  fork  usa  esse  endereço  diretamente  como  
um  endereço  virtual  ao  copiar  a  memória  do  usuário  do  processo  pai  para  o  filho.
 Há  alguns  endereços  virtuais  do  kernel  que  não  são  mapeados  diretamente:
 •  A  página  do  trampolim.  Ela  é  mapeada  no  topo  do  espaço  de  endereço  virtual;  as  tabelas  de  páginas  do  usuário  
têm  esse  mesmo  mapeamento.  O  Capítulo  4  discute  o  papel  da  página  do  trampolim,  mas  vemos  aqui  um  caso  de  
uso  interessante  de  tabelas  de  páginas;  uma  página  física  (contendo  o  código  do  trampolim)  é  mapeada  duas  
vezes  no  espaço  de  endereço  virtual  do  kernel:  uma  vez  no  topo  do  espaço  de  endereço  virtual  e  outra  com  um  
mapeamento  direto.
 •  As  páginas  da  pilha  do  kernel.  Cada  processo  tem  sua  própria  pilha  do  kernel,  que  é  mapeada  em  uma  posição  alta  
para  que,  abaixo  dela,  o  xv6  possa  deixar  uma  página  de  guarda  não  mapeada.  A  PTE  da  página  de  guarda  é  
inválida  (ou  seja,  PTE_V  não  está  definida),  de  modo  que,  se  o  kernel  transbordar  uma  pilha  do  kernel,  
provavelmente  causará  uma  exceção  e  o  kernel  entrará  em  pânico.  Sem  uma  página  de  guarda,  uma  pilha  
transbordando  sobrescreveria  outra  memória  do  kernel,  resultando  em  operação  incorreta.  Uma  falha  de  pânico  é  preferível.
 Embora  o  kernel  utilize  suas  pilhas  por  meio  de  mapeamentos  de  alta  memória,  elas  também  são  acessíveis  ao  
kernel  por  meio  de  um  endereço  mapeado  diretamente.  Um  projeto  alternativo  poderia  ter  apenas  o  mapeamento  direto  
e  usar  as  pilhas  no  endereço  mapeado  diretamente.  Nesse  arranjo,  no  entanto,  fornecer  páginas  de  guarda  envolveria  o  
desmapeamento  de  endereços  virtuais  que,  de  outra  forma,  se  refeririam  à  memória  física,  o  que  dificultaria  o  uso.
 O  kernel  mapeia  as  páginas  do  trampolim  e  o  texto  do  kernel  com  as  permissões  PTE_R  e  PTE_X.  O  kernel  lê  e  
executa  instruções  dessas  páginas.  O  kernel  mapeia  as  outras  páginas  com  as  permissões  PTE_R  e  PTE_W,  para  que  
possa  ler  e  gravar  na  memória  dessas  páginas.  Os  mapeamentos  para  as  páginas  de  guarda  são  inválidos.
 3.3  Código:  criando  um  espaço  de  endereço
 A  maior  parte  do  código  xv6  para  manipular  espaços  de  endereço  e  tabelas  de  páginas  reside  em  vm.c  ( kernel/vm.c:1).  
A  estrutura  de  dados  central  é  pagetable_t,  que  na  verdade  é  um  ponteiro  para  um  RISC-V
 35
Machine Translated by Google
 Tabela  de  páginas  raiz;  uma  pagetable_t  pode  ser  a  tabela  de  páginas  do  kernel  ou  uma  das  tabelas  de  páginas  por  
processo.  As  funções  centrais  são  walk,  que  encontra  o  PTE  para  um  endereço  virtual,  e  mappages,  que  instala  PTEs  
para  novos  mapeamentos.  Funções  que  começam  com  kvm  manipulam  a  tabela  de  páginas  do  kernel;  funções  que  
começam  com  uvm  manipulam  uma  tabela  de  páginas  do  usuário;  outras  funções  são  usadas  para  ambos.  copyout  e  
copyin  copiam  dados  de  e  para  endereços  virtuais  do  usuário  fornecidos  como  argumentos  de  chamada  de  sistema;  elas  
estão  em  vm.c  porque  precisam  traduzir  explicitamente  esses  endereços  para  encontrar  a  memória  física  correspondente.
 No  início  da  sequência  de  inicialização,  o  principal  chama  kvminit  (kernel/vm.c:54)  para  criar  a  tabela  de  páginas  do  
kernel  usando  kvmmake  (kernel/vm.c:20).  Esta  chamada  ocorre  antes  que  o  xv6  habilite  a  paginação  no  RISC-V,  
portanto,  os  endereços  se  referem  diretamente  à  memória  física.  O  kvmmake  primeiro  aloca  uma  página  de  memória  
física  para  armazenar  a  página  da  tabela  de  páginas  raiz.  Em  seguida,  ele  chama  o  kvmmap  para  instalar  as  traduções  
necessárias  ao  kernel.  As  traduções  incluem  as  instruções  e  os  dados  do  kernel,  a  memória  física  até  o  PHYSTOP  e  os  
intervalos  de  memória  que,  na  verdade,  são  dispositivos.  proc_mapstacks  (kernel/proc.c:33)  aloca  uma  pilha  do  kernel  
para  cada  processo.  Ele  chama  o  kvmmap  para  mapear  cada  pilha  no  endereço  virtual  gerado  pelo  KSTACK,  o  que  
deixa  espaço  para  as  páginas  de  proteção  de  pilha  inválidas.
 kvmmap  (kernel/vm.c:132)  chama  mappages  (kernel/vm.c:143),  que  instala  mapeamentos  em  uma  tabela  de  
páginas  para  um  intervalo  de  endereços  virtuais  em  um  intervalo  correspondente  de  endereços  físicos.  Ele  faz  isso  
separadamente  para  cada  endereço  virtual  no  intervalo,  em  intervalos  de  página.  Para  cada  endereço  virtual  a  ser  
mapeado,  mappages  chama  walk  para  encontrar  o  endereço  da  PTE  correspondente  a  esse  endereço.  Em  
seguida,  ele  inicializa  a  PTE  para  conter  o  número  da  página  física  relevante,  as  permissões  desejadas  (PTE_W,  
PTE_X  e/ou  PTE_R)  e  PTE_V  para  marcar  a  PTE  como  válida  (kernel/vm.c:158).
 andar  (kernel/vm.c:86)  imita  o  hardware  de  paginação  RISC-V  enquanto  procura  no  PTE  por  um  endereço  virtual  
(veja  a  Figura  3.2).  O  comando  walk  desce  a  tabela  de  páginas  de  3  níveis  9  bits  por  vez.  Ele  usa  os  9  bits  de  endereço  
virtual  de  cada  nível  para  encontrar  o  PTE  da  tabela  de  páginas  do  próximo  nível  ou  da  página  final  (kernel/vm.c:92).  Se  
o  PTE  não  for  válido,  a  página  necessária  ainda  não  foi  alocada;  se  o  argumento  alloc  estiver  definido,  walk  aloca  uma  
nova  página  da  tabela  de  páginas  e  insere  seu  endereço  físico  no  PTE.  Ele  retorna  o  endereço  do  PTE  na  camada  mais  
baixa  da  árvore  (kernel/vm.c:102).
 O  código  acima  depende  da  memória  física  ser  mapeada  diretamente  no  espaço  de  endereço  virtual  do  kernel.  Por  
exemplo,  à  medida  que  o  walk  desce  níveis  da  tabela  de  páginas,  ele  extrai  o  endereço  (físico)  da  tabela  de  páginas  do  
nível  seguinte  de  uma  PTE  (kernel/vm.c:94).  e  então  usa  esse  endereço  como  um  endereço  virtual  para  buscar  o  PTE  
no  próximo  nível  abaixo  (kernel/vm.c:92).
 chamadas  principais  kvminithart  (kernel/vm.c:62)  para  instalar  a  tabela  de  páginas  do  kernel.  Ele  grava  o  
endereço  físico  da  página  raiz  da  tabela  de  páginas  no  registrador  satp.  Depois  disso,  a  CPU  traduzirá  os  
endereços  usando  a  tabela  de  páginas  do  kernel.  Como  o  kernel  usa  um  mapeamento  de  identidade,  o  endereço  
virtual  da  próxima  instrução  será  mapeado  para  o  endereço  de  memória  física  correto.
 Cada  CPU  RISC-V  armazena  em  cache  as  entradas  da  tabela  de  páginas  em  um  Translation  Look-aside  Buffer  
(TLB)  e,  quando  o  xv6  altera  uma  tabela  de  páginas,  deve  informar  à  CPU  para  invalidar  as  entradas  TLB  em  cache  
correspondentes.  Caso  contrário,  em  algum  momento  posterior,  o  TLB  poderá  usar  um  mapeamento  em  cache  antigo,  
apontando  para  uma  página  física  que,  nesse  meio  tempo,  foi  alocada  a  outro  processo  e,  como  resultado,  um  processo  
poderá  rabiscar  na  memória  de  outro  processo.  O  RISC-V  possui  uma  instrução  sfence.vma  que  limpa  o  TLB  da  CPU  
atual.  O  xv6  executa  sfence.vma  em  kvminithart  após  recarregar  o  registrador  satp  e  no  código  trampoline  que  alterna  
para  um
 36
Machine Translated by Google
 tabela  de  página  do  usuário  antes  de  retornar  ao  espaço  do  usuário  (kernel/trampoline.S:89).
 Também  é  necessário  emitir  sfence.vma  antes  de  alterar  o  satp,  para  aguardar  a  conclusão  de  todos  os  carregamentos  e  
armazenamentos  pendentes.  Essa  espera  garante  que  as  atualizações  anteriores  na  tabela  de  páginas  tenham  sido  concluídas  e  
que  os  carregamentos  e  armazenamentos  anteriores  usem  a  tabela  de  páginas  antiga,  não  a  nova.
 um.
 Para  evitar  a  limpeza  completa  do  TLB,  as  CPUs  RISC-V  podem  suportar  identificadores  de  espaço  de  endereço  
(ASIDs)  [3].  O  kernel  pode  então  limpar  apenas  as  entradas  TLB  para  um  espaço  de  endereço  específico.  O  Xv6  não  
utiliza  esse  recurso.
 3.4  Alocação  de  memória  física
 O  kernel  deve  alocar  e  liberar  memória  física  em  tempo  de  execução  para  tabelas  de  páginas,  memória  do  usuário,  pilhas  do  kernel  
e  buffers  de  pipe.
 O  Xv6  utiliza  a  memória  física  entre  o  final  do  kernel  e  o  PHYSTOP  para  alocação  em  tempo  de  execução.  Ele  aloca  e  libera  
páginas  inteiras  de  4.096  bytes  por  vez.  Ele  rastreia  quais  páginas  estão  livres,  encadeando  uma  lista  encadeada  entre  as  próprias  
páginas.  A  alocação  consiste  em  remover  uma  página  da  lista  encadeada;  a  liberação  consiste  em  adicionar  a  página  liberada  à  
lista.
 3.5  Código:  Alocador  de  memória  física
 O  alocador  reside  em  kalloc.c  (kernel/kalloc.c:1).  A  estrutura  de  dados  do  alocador  é  uma  lista  livre  de  páginas  de  memória  física  
disponíveis  para  alocação.  Cada  elemento  da  lista  de  páginas  livres  é  uma  execução  de  struct  (kernel/kalloc.c:17).  De  onde  o  
alocador  obtém  a  memória  para  armazenar  essa  estrutura  de  dados?  Ele  armazena  a  estrutura  de  execução  de  cada  página  livre  
na  própria  página  livre,  já  que  não  há  mais  nada  armazenado  lá.  A  lista  livre  é  protegida  por  um  bloqueio  de  spin  (kernel/
 kalloc.c:21-24).  A  lista  e  o  bloqueio  são  encapsulados  em  uma  struct  para  deixar  claro  que  o  bloqueio  protege  os  campos  na  struct.  
Por  enquanto,  ignore  o  bloqueio  e  as  chamadas  para  acquire  e  release;  o  Capítulo  6  examinará  o  bloqueio  em  detalhes.
 A  função  main  chama  kinit  para  inicializar  o  alocador  (kernel/kalloc.c:27).  O  kinit  inicializa  a  lista  de  páginas  livres  para  
armazenar  todas  as  páginas  entre  o  final  do  kernel  e  o  PHYSTOP.  O  Xv6  deve  determinar  quanta  memória  física  está  disponível  
analisando  as  informações  de  configuração  fornecidas  pelo  hardware.  Em  vez  disso,  o  xv6  assume  que  a  máquina  possui  128  
megabytes  de  RAM.  O  kinit  chama  o  freerange  para  adicionar  memória  à  lista  de  páginas  livres  por  meio  de  chamadas  por  página  
ao  kfree.  Uma  PTE  só  pode  se  referir  a  um  endereço  físico  alinhado  em  um  limite  de  4096  bytes  (é  um  múltiplo  de  4096),  portanto,  
o  freerange  usa  PGROUNDUP  para  garantir  que  libere  apenas  os  endereços  físicos  alinhados.  O  alocador  inicia  sem  memória;  
essas  chamadas  ao  kfree  lhe  dão  algum  trabalho  para  gerenciar.
 O  alocador  às  vezes  trata  endereços  como  inteiros  para  realizar  operações  aritméticas  com  eles  (por  exemplo,  percorrendo  
todas  as  páginas  em  modo  livre)  e,  às  vezes,  usa  endereços  como  ponteiros  para  ler  e  escrever  na  memória  (por  exemplo,  
manipulando  a  estrutura  de  execução  armazenada  em  cada  página);  esse  uso  duplo  de  endereços  é  o  principal  motivo  pelo  qual  o  
código  do  alocador  está  cheio  de  conversões  de  tipo  em  C.  O  outro  motivo  é  que  a  liberação  e  a  alocação  alteram  inerentemente  o  
tipo  de  memória.
 37
Machine Translated by Google
 A  função  kfree  (kernel/kalloc.c:47)  começa  definindo  cada  byte  na  memória  que  está  sendo  liberada  para  o  valor  1.  Isso  fará  
com  que  o  código  que  usa  memória  após  liberá-la  (usa  “referências  pendentes”)  leia  lixo  em  vez  do  antigo  conteúdo  válido;  
esperançosamente,  isso  fará  com  que  esse  código  quebre  mais  rápido.
 Em  seguida,  kfree  anexa  a  página  à  lista  livre:  ele  converte  pa  para  um  ponteiro  para  struct  run,  registra  o  início  antigo  
da  lista  livre  em  r->next  e  define  a  lista  livre  como  igual  a  r.  kalloc  remove  e  retorna  o  primeiro  elemento  na  lista  livre.
 3.6  Espaço  de  endereço  do  processo
 Cada  processo  possui  uma  tabela  de  páginas  separada  e,  quando  o  xv6  alterna  entre  processos,  também  altera  as  
tabelas  de  páginas.  A  Figura  3.4  mostra  o  espaço  de  endereço  de  um  processo  com  mais  detalhes  do  que  a  Figura  2.3.  
A  memória  de  usuário  de  um  processo  começa  no  endereço  virtual  zero  e  pode  crescer  até  MAXVA  (kernel/riscv.h:360).  
permitindo  que  um  processo  enderece  em  princípio  256  Gigabytes  de  memória.
 O  espaço  de  endereço  de  um  processo  consiste  em  páginas  que  contêm  o  texto  do  programa  (que  o  xv6  mapeia  
com  as  permissões  PTE_R,  PTE_X  e  PTE_U),  páginas  que  contêm  os  dados  pré-inicializados  do  programa,  uma  página  
para  a  pilha  e  páginas  para  o  heap.  O  xv6  mapeia  os  dados,  a  pilha  e  o  heap  com  as  permissões  PTE_R,  PTE_W  e  
PTE_U.
 Usar  permissões  dentro  de  um  espaço  de  endereço  de  usuário  é  uma  técnica  comum  para  proteger  um  processo  de  usuário.
 Se  o  texto  fosse  mapeado  com  PTE_W,  um  processo  poderia  modificar  acidentalmente  seu  próprio  programa;  por  exemplo,  um  
erro  de  programação  pode  fazer  com  que  o  programa  escreva  em  um  ponteiro  nulo,  modificando  instruções  no  endereço  0,  e  
então  continue  em  execução,  talvez  causando  mais  confusão.  Para  detectar  esses  erros  imediatamente,  o  xv6  mapeia  o  texto  
sem  PTE_W;  se  um  programa  tentar  acidentalmente  armazenar  no  endereço  0,  o  hardware  se  recusará  a  executar  o  
armazenamento  e  gerará  uma  falha  de  página  (consulte  a  Seção  4.6).
 O  kernel  então  mata  o  processo  e  imprime  uma  mensagem  informativa  para  que  o  desenvolvedor  possa  rastrear  o  problema.
 Da  mesma  forma,  ao  mapear  dados  sem  PTE_X,  um  programa  de  usuário  não  pode  pular  acidentalmente  para  um
 endereço  nos  dados  do  programa  e  começar  a  executar  nesse  endereço.
 No  mundo  real,  fortalecer  um  processo  por  meio  da  definição  cuidadosa  de  permissões  também  auxilia  na  defesa  contra  
ataques  de  segurança.  Um  adversário  pode  fornecer  uma  entrada  cuidadosamente  construída  a  um  programa  (por  exemplo,  um  
servidor  web)  que  acione  um  bug  no  programa  na  esperança  de  transformá-lo  em  um  exploit  [14].
 Definir  permissões  cuidadosamente  e  outras  técnicas,  como  a  randomização  do  layout  do  espaço  de  endereço  do  usuário,  tornam  
esses  ataques  mais  difíceis.
 A  pilha  é  uma  única  página  e  é  exibida  com  o  conteúdo  inicial  criado  por  exec.  Strings  contendo  os  argumentos  da  linha  de  
comando,  bem  como  um  array  de  ponteiros  para  eles,  ficam  no  topo  da  pilha.  Logo  abaixo,  estão  os  valores  que  permitem  que  
um  programa  inicie  em  main  como  se  a  função  main(argc,  argv)  tivesse  acabado  de  ser  chamada.
 Para  detectar  que  uma  pilha  de  usuário  está  transbordando  a  memória  alocada,  o  xv6  coloca  uma  página  de  guarda  
inacessível  logo  abaixo  da  pilha,  limpando  o  sinalizador  PTE_U .  Se  a  pilha  de  usuário  transbordar  e  o  processo  tentar  usar  um  
endereço  abaixo  da  pilha,  o  hardware  gerará  uma  exceção  de  falha  de  página  porque  a  página  de  guarda  está  inacessível  para  
um  programa  em  execução  no  modo  de  usuário.  Um  sistema  operacional  real  pode,  em  vez  disso,  alocar  automaticamente  mais  
memória  para  a  pilha  de  usuário  quando  ela  transborda.
 Quando  um  processo  solicita  ao  xv6  mais  memória  do  usuário,  o  xv6  aumenta  o  heap  do  processo.  O  xv6  usa  primeiro
 38
Machine Translated by Google
 MAXVA
 TAMANHO  DA  PÁGINA
 Alinhamento  de  página
 0
 Receita--
estrutura  de  
armadilha  de  trampolim
 não  utilizado
 monte
 pilha
 pilha  
página  de  guarda
 dados
 não  utilizado
 texto
 RW-
R-WU
R-WU
R-WU
R-XU
 argumento  0
 ...
 argumento  N  0
 endereço  do  argumento  N
 ...
 endereço  do  argumento  0  
endereço  do  endereço  do  
argumento  0  
argc  
0xFFFFFFF
 (vazio)
 string  terminada  em  nulo  argv[argc]
 argv[0]
 argumento  argv  do  principal
 argumento  argc  do  retorno  
principal  PC  para  principal
 Figura  3.4:  Espaço  de  endereço  de  usuário  de  um  processo,  com  sua  pilha  inicial.
 kalloc  para  alocar  páginas  físicas.  Em  seguida,  ele  adiciona  PTEs  à  tabela  de  páginas  do  processo  que  apontam  para  
as  novas  páginas  físicas.  O  Xv6  define  os  sinalizadores  PTE_W,  PTE_R,  PTE_U  e  PTE_V  nesses  PTEs.  A  maioria  
dos  processos  não  utiliza  todo  o  espaço  de  endereço  do  usuário;  o  Xv6  deixa  PTE_V  livre  em  PTEs  não  utilizados.
 Vemos  aqui  alguns  bons  exemplos  do  uso  de  tabelas  de  páginas.  Primeiro,  as  tabelas  de  páginas  de  diferentes  
processos  traduzem  endereços  de  usuários  para  diferentes  páginas  de  memória  física,  de  modo  que  cada  processo  tenha  
memória  de  usuário  privada.  Segundo,  cada  processo  vê  sua  memória  como  tendo  endereços  virtuais  contíguos  começando  
em  zero,  enquanto  a  memória  física  do  processo  pode  ser  não  contígua.  Terceiro,  o  kernel  mapeia  uma  página  com  código  
trampolim  no  topo  do  espaço  de  endereço  do  usuário  (sem  PTE_U),  portanto,  uma  única  página  de  memória  física  aparece  
em  todos  os  espaços  de  endereço,  mas  pode  ser  usada  apenas  pelo  kernel.
 3.7  Código:  sbrk
 sbrk  é  a  chamada  de  sistema  para  que  um  processo  reduza  ou  aumente  sua  memória.  A  chamada  de  sistema  é  
implementada  pela  função  growproc  (kernel/proc.c:260).  growproc  chama  uvmalloc  ou  uvmdealloc,  dependendo  se  n  
é  positivo  ou  negativo.  uvmalloc  (kernel/vm.c:226)  aloca  memória  física  com  kalloc  e  adiciona  PTEs  à  tabela  de  
páginas  do  usuário  com  mappages.  uvmdealloc  chama  uvmunmap  (kernel/vm.c:171),  que  usa  walk  para  encontrar  
PTEs  e  kfree  para  liberar  a  memória  física  à  qual  se  referem.
 O  Xv6  usa  a  tabela  de  páginas  de  um  processo  não  apenas  para  informar  ao  hardware  como  mapear  os  endereços  
virtuais  do  usuário,  mas  também  como  o  único  registro  de  quais  páginas  de  memória  física  são  alocadas  para  aquele  
processo.  É  por  isso  que  a  liberação  de  memória  do  usuário  (no  uvmunmap)  requer  a  análise  da  tabela  de  páginas  do  usuário.
 39
Machine Translated by Google
 3.8  Código:  exec
 exec  é  uma  chamada  de  sistema  que  substitui  o  espaço  de  endereço  do  usuário  de  um  processo  por  dados  lidos  de  um  
arquivo,  chamado  de  arquivo  binário  ou  executável.  Um  binário  normalmente  é  a  saída  do  compilador  e  do  vinculador  e  
contém  instruções  de  máquina  e  dados  do  programa.  exec  (kernel/exec.c:23)  abre  o  caminho  binário  nomeado  usando  
namei  (kernel/exec.c:36),  que  é  explicado  no  Capítulo  8.  Em  seguida,  ele  lê  o  cabeçalho  ELF.  Os  binários  Xv6  são  
formatados  no  formato  ELF  amplamente  utilizado,  definido  em  (kernel/elf.h).  Um  binário  ELF  consiste  em  um  cabeçalho  
ELF,  struct  elfhdr  (kernel/elf.h:6),  seguido  por  uma  sequência  de  cabeçalhos  de  seção  do  programa,  struct  proghdr  (kernel/
 elf.h:25).  Cada  progvhdr  descreve  uma  seção  do  aplicativo  que  deve  ser  carregada  na  memória;  os  programas  xv6  têm  
dois  cabeçalhos  de  seção  de  programa:  um  para  instruções  e  um  para  dados.
 O  primeiro  passo  é  uma  verificação  rápida  para  verificar  se  o  arquivo  provavelmente  contém  um  binário  ELF.  Um  
binário  ELF  começa  com  o  "número  mágico"  de  quatro  bytes  0x7F,  'E',  'L',  'F'  ou  ELF_MAGIC  (kernel/elf.h:3).  Se  o  
cabeçalho  ELF  tiver  o  número  mágico  correto,  exec  assume  que  o  binário  está  bem  formado.
 exec  aloca  uma  nova  tabela  de  páginas  sem  mapeamentos  de  usuários  com  proc_pagetable  (kernel/exec.c:49),
 aloca  memória  para  cada  segmento  ELF  com  uvmalloc  (kernel/exec.c:65),  e  carrega  cada  segmento  na  memória  com  
loadseg  (kernel/exec.c:10).  loadseg  usa  walkaddr  para  encontrar  o  endereço  físico  da  memória  alocada  na  qual  gravar  
cada  página  do  segmento  ELF  e  readi  para  ler  o  arquivo.
 O  cabeçalho  da  seção  do  programa  para /init,  o  primeiro  programa  de  usuário  criado  com  exec,  se  parece  com  isto:
 #  objdump  -p  usuário/_init
 usuário/_init:
 Cabeçalho  do  programa:  
0x70000003  off
 CARREGAR  desligado
 formato  de  arquivo  elf64-little
 0x00000000000006bb0  vaddr  0x0000000000000000
 paddr  0x00000000000000000  align  2**0  filesz  0x000000000000004a  
memsz  0x0000000000000000  sinalizadores  r-
0x0000000000001000  vaddr  0x0000000000000000
 paddr  0x00000000000000000  alinhar  2**12  arquivosz  0x0000000000001000  
memsz  0x000000000001000  sinalizadores  rx
 CARREGAR  desligado
 0x0000000000002000  vaddr  0x000000000001000
 memsz  0x000000000000030  sinalizadores  rw
EMPILHAMENTO  desligado
 paddr  0x0000000000001000  alinhar  2**12  arquivosz  0x0000000000000010  
0x00000000000000000  vaddr  0x000000000000000  paddr  0x0000000000000000  
alinhar  2**4  arquivosz  0x0000000000000000  memsz  
0x0000000000000000  sinalizadores  rw
Vemos  que  o  segmento  de  texto  deve  ser  carregado  no  endereço  virtual  0x0  na  memória  (sem  permissões  de  
gravação)  a  partir  do  conteúdo  no  deslocamento  0x1000  do  arquivo.  Também  vemos  que  os  dados  devem  ser  carregados  
no  endereço  0x1000,  que  está  no  limite  de  uma  página,  e  sem  permissões  de  execução.
 O  valor  "filesz"  do  cabeçalho  de  uma  seção  de  programa  pode  ser  menor  que  o  valor  "memsz",  indicando  que  o  
espaço  entre  eles  deve  ser  preenchido  com  zeros  (para  variáveis  globais  em  C)  em  vez  de  lido  do  arquivo.  Para /init,  o  
valor  "filesz"  de  dados  é  0x10  bytes  e  o  valor  "memsz"  é  0x30  bytes,  e,  portanto,  uvmalloc  aloca  memória  física  suficiente  
para  armazenar  0x30  bytes,  mas  lê  apenas  0x10  bytes  do  arquivo /init.
 40
Machine Translated by Google
 Agora,  exec  aloca  e  inicializa  a  pilha  do  usuário.  Ele  aloca  apenas  uma  página  de  pilha.  exec  copia  as  strings  de  
argumentos  para  o  topo  da  pilha,  uma  de  cada  vez,  registrando  os  ponteiros  para  elas  em  ustack.
 Ele  coloca  um  ponteiro  nulo  no  final  do  que  será  a  lista  argv  passada  para  main.  As  três  primeiras  entradas  em  ustack  são  
o  contador  de  programa  de  retorno  falso,  argc  e  o  ponteiro  argv .
 O  comando  exec  coloca  uma  página  inacessível  logo  abaixo  da  página  da  pilha,  de  modo  que  programas  que  tentarem  
usar  mais  de  uma  página  apresentarão  falhas.  Essa  página  inacessível  também  permite  que  o  comando  exec  trate  de  
argumentos  muito  grandes;  nessa  situação,  a  cópia  (kernel/vm.c:352)  a  função  que  exec  usa  para  copiar  argumentos  para  
a  pilha  notará  que  a  página  de  destino  não  está  acessível  e  retornará  -1.
 Durante  a  preparação  da  nova  imagem  de  memória,  se  exec  detectar  um  erro,  como  um  segmento  de  programa  
inválido,  ele  salta  para  o  rótulo  "bad",  libera  a  nova  imagem  e  retorna  -1.  exec  deve  aguardar  para  liberar  a  imagem  antiga  
até  ter  certeza  de  que  a  chamada  de  sistema  será  bem-sucedida:  se  a  imagem  antiga  desaparecer,  a  chamada  de  sistema  
não  poderá  retornar  -1  para  ela.  Os  únicos  casos  de  erro  em  exec  ocorrem  durante  a  criação  da  imagem.  Assim  que  a  
imagem  estiver  concluída,  exec  pode  confirmar  a  nova  tabela  de  páginas  (kernel/exec.c:125).  e  liberar  o  antigo  (kernel/
 exec.c:129).
 exec  carrega  bytes  do  arquivo  ELF  para  a  memória  em  endereços  especificados  pelo  arquivo  ELF.  Usuários  ou  
processos  podem  colocar  quaisquer  endereços  que  desejarem  em  um  arquivo  ELF.  Portanto,  exec  é  arriscado,  pois  os  
endereços  no  arquivo  ELF  podem  se  referir  ao  kernel,  acidentalmente  ou  propositalmente.  As  consequências  para  um  
kernel  desavisado  podem  variar  de  uma  falha  a  uma  subversão  maliciosa  dos  mecanismos  de  isolamento  do  kernel  (ou  
seja,  uma  falha  de  segurança).  O  Xv6  realiza  uma  série  de  verificações  para  evitar  esses  riscos.  Por  exemplo,  if(ph.vaddr  
+  ph.memsz  <  ph.vaddr)  verifica  se  a  soma  excede  um  inteiro  de  64  bits.  O  perigo  é  que  um  usuário  possa  construir  um  
binário  ELF  com  um  ph.vaddr  que  aponte  para  um  endereço  escolhido  pelo  usuário  e  um  ph.memsz  grande  o  suficiente  
para  que  a  soma  exceda  para  0x1000,  o  que  parecerá  um  valor  válido.  Em  uma  versão  mais  antiga  do  xv6,  na  qual  o  
espaço  de  endereço  do  usuário  também  continha  o  kernel  (mas  não  era  legível/gravável  no  modo  usuário),  o  usuário  podia  
escolher  um  endereço  que  correspondesse  à  memória  do  kernel  e,  assim,  copiar  os  dados  do  binário  ELF  para  o  kernel.  
Na  versão  RISC-V  do  xv6,  isso  não  acontece,  pois  o  kernel  tem  sua  própria  tabela  de  páginas;  loadseg  carrega  na  tabela  
de  páginas  do  processo,  não  na  tabela  de  páginas  do  kernel.
 É  fácil  para  um  desenvolvedor  de  kernel  omitir  uma  verificação  crucial,  e  kernels  do  mundo  real  têm  um  longo  histórico  
de  verificações  ausentes,  cuja  ausência  pode  ser  explorada  por  programas  de  usuário  para  obter  privilégios  de  kernel.  É  
provável  que  o  xv6  não  faça  um  trabalho  completo  de  validação  dos  dados  de  nível  de  usuário  fornecidos  ao  kernel,  o  que  
um  programa  de  usuário  malicioso  pode  explorar  para  contornar  o  isolamento  do  xv6.
 3.9  Mundo  real
 Como  a  maioria  dos  sistemas  operacionais,  o  xv6  usa  o  hardware  de  paginação  para  proteção  e  mapeamento  de  memória.
 A  maioria  dos  sistemas  operacionais  faz  uso  muito  mais  sofisticado  de  paginação  do  que  o  xv6,  combinando  exceções  de  
paginação  e  falha  de  página,  que  discutiremos  no  Capítulo  4.
 O  Xv6  é  simplificado  pelo  uso,  pelo  kernel,  de  um  mapeamento  direto  entre  endereços  virtuais  e  físicos,  e  pela  
suposição  de  que  há  RAM  física  no  endereço  0x8000000,  onde  o  kernel  espera  ser  carregado.  Isso  funciona  com  o  QEMU,  
mas  em  hardware  real  acaba  sendo  uma  má  ideia;  o  hardware  real  coloca  RAM  e  dispositivos  em  endereços  físicos  
imprevisíveis,  de  modo  que  (por  exemplo)  pode  não  haver  RAM  em  0x8000000,  onde  o  xv6  espera  poder  armazenar  o  
kernel.  Um  kernel  mais  sério
 41
Machine Translated by Google
 os  projetos  exploram  a  tabela  de  páginas  para  transformar  layouts  arbitrários  de  memória  física  de  hardware  em  layouts  previsíveis  
de  endereços  virtuais  do  kernel.
 O  RISC-V  suporta  proteção  no  nível  de  endereços  físicos,  mas  o  xv6  não  usa  esse  recurso.
 Em  máquinas  com  muita  memória,  pode  fazer  sentido  usar  o  suporte  do  RISC-V  para  "superpáginas".  Páginas  pequenas  
fazem  sentido  quando  a  memória  física  é  pequena,  para  permitir  alocação  e  paginação  para  o  disco  com  granularidade  fina.  Por  
exemplo,  se  um  programa  usa  apenas  8  quilobytes  de  memória,  atribuir  a  ele  uma  superpágina  inteira  de  4  megabytes  de  memória  
física  é  um  desperdício.  Páginas  maiores  fazem  sentido  em  máquinas  com  muita  RAM  e  podem  reduzir  a  sobrecarga  na  manipulação  
da  tabela  de  páginas.
 A  falta  de  um  alocador  do  tipo  malloc  no  kernel  xv6  que  possa  fornecer  memória  para  objetos  pequenos  impede  que  o  kernel  
use  estruturas  de  dados  sofisticadas  que  exigiriam  alocação  dinâmica.
 Um  kernel  mais  elaborado  provavelmente  alocaria  muitos  tamanhos  diferentes  de  blocos  pequenos,  em  vez  de  (como  no  xv6)  
apenas  blocos  de  4.096  bytes;  um  alocador  de  kernel  real  precisaria  lidar  com  alocações  pequenas  e  também  grandes.
 A  alocação  de  memória  é  um  tópico  recorrente  e  polêmico,  cujos  problemas  básicos  são  o  uso  eficiente  de  memória  limitada  e  
a  preparação  para  solicitações  futuras  desconhecidas  [9].  Hoje  em  dia,  as  pessoas  se  preocupam  mais  com  velocidade  do  que  com  
eficiência  de  espaço.
